\section{第一步：固定点}

几何作图的第一步是定义固定点，图形将从这些点构建。

总体思路是避免操作坐标，而是使用第一步中固定的点或通过宏包工具获得的点的引用。即使可以直接使用坐标，我认为这也是个坏主意：最好使用命名点。

\tkzname{\tkznameofpack} 使用专用于几何作图的宏和词汇。你当然可以使用 \TIKZ 的工具，但在我看来不混合这两种语法更合乎逻辑。

对于 \TIKZ，\tkzname{\tkznameofpack} 的点是一个特殊的 \code{node}。在下一节中我们将看到如何用坐标定义点。这里不讨论点的样式（颜色和形状）：你会在示例中找到一些提示，并在第 \ref{custom} 节中找到更多信息。


\section{点的定义：\tkzcname{tkzDefPoint} 或 \tkzcname{tkzDefPoints}}

 点可以通过以下任一方式指定：
\begin{itemize}
\item 笛卡尔坐标；
\item 极坐标；
\item 命名点；
\item 相对点。
\end{itemize}


当一个点的名称与唯一一对十进制数关联时，该点就被定义：
 $(x,y)$（$x$ 横坐标，$y$ 纵坐标）或 $(a:d)$（$a$ 角，$d$ 距离）。
 这之所以可能，是因为平面被赋予了一个单位为 $1$~cm 的标准正交笛卡尔坐标系。

 笛卡尔坐标 $(a,b)$ 指向在 $x$ 方向上 $a$ 厘米、在 $y$ 方向上 $b$ 厘米的点。

 极坐标中的一个点需要一个以度为单位的角 $\alpha$，以及到原点的距离 $d$；默认单位是 \texttt{cm}。

 宏 \tkzNameMacro{tkzDefPoint} 通过给点赋值坐标来定义点。这个宏基于

 \tkzNameMacro{coordinate}，这是 \TIKZ 的一个宏，并接受 \TIKZ 选项如 \tkzname{shift}。计算由宏包 \tkzNamePack{xfp} 完成。我们可以使用笛卡尔坐标或极坐标。

\begin{minipage}[t]{0.48\textwidth}
 笛卡尔坐标
\begin{tkzexample}[code only,small]
\begin{tikzpicture}[scale=1] 
  \tkzInit[xmax=5,ymax=5]
  % necessary to limit
  % the size of the axes
  \tkzDrawX[>=latex]
  \tkzDrawY[>=latex]
  \tkzDefPoints{0/0/O,1/0/I,0/1/J}
  \tkzDefPoint(3,4){A} 
  \tkzDrawPoints(O,A) 
  \tkzLabelPoint[above](A){$A_1(x_1,y_1)$} 
  \tkzShowPointCoord[xlabel=$x_1$,
                     ylabel=$y_1$](A) 
  \tkzLabelPoints(O,I)
  \tkzLabelPoints[left](J)
  \tkzDrawPoints[shape=cross](I,J) 
\end{tikzpicture}
\end{tkzexample}%
\end{minipage}
\begin{minipage}[t]{0.45\textwidth}
 极坐标
\begin{tkzexample}[code only,small]
\begin{tikzpicture}[scale=1]
  \tkzInit[xmax=5,ymax=5]
  \tkzDrawX[>=latex]
  \tkzDrawY[>=latex]  
  \tkzDefPoints{0/0/O,1/0/I,0/1/J}
  \tkzDefPoint(40:4){P} 
  \tkzDrawSegment[dim={$d$,
                 16pt,above=6pt}](O,P)
  \tkzDrawPoints(O,P) 
  \tkzMarkAngle[mark=none,->](I,O,P) 
  \tkzFillAngle[opacity=.5](I,O,P) 
  \tkzLabelAngle[pos=1.25](I,O,P){%
                              $\alpha$}
  \tkzLabelPoint[right](P){$P(\alpha:d)$}
  \tkzDrawPoints[shape=cross](I,J)
  \tkzLabelPoints(O,I)
  \tkzLabelPoints[left](J)
\end{tikzpicture}
\end{tkzexample}
\end{minipage}%

\begin{minipage}[b]{0.45\textwidth}
\begin{tikzpicture}[scale=1]
  \tkzInit[xmax=5,ymax=5]
  \tkzDrawX[>=latex]
  \tkzDrawY[>=latex]
  \tkzDefPoints{0/0/O,1/0/I,0/1/J}
  \tkzDefPoint(3,4){A}
  \tkzDrawPoints(O,A)
  \tkzLabelPoint[above](A){$A_1 (x_1,y_1)$}
  \tkzShowPointCoord[xlabel=$x_1$,ylabel=$y_1$](A)
  \tkzLabelPoints(O,I)
  \tkzLabelPoints[left](J)
  \tkzDrawPoints[shape=cross](I,J)
\end{tikzpicture}
\end{minipage}
\begin{minipage}[b]{0.45\textwidth}
\begin{tikzpicture}[,scale=1]
  \tkzInit[xmax=5,ymax=5]
  \tkzDrawX[>=latex]
  \tkzDrawY[>=latex]  
  \tkzDefPoints{0/0/O,1/0/I,0/1/J}
  \tkzDefPoint(40:4){P} 
  \tkzDrawSegment[dim={$d$,
                 16pt,above=6pt}](O,P)
  \tkzDrawPoints(O,P) 
  \tkzMarkAngle[mark=none,->](I,O,P) 
  \tkzFillAngle[opacity=.5](I,O,P) 
  \tkzLabelAngle[pos=1.25](I,O,P){$\alpha$}  
  \tkzLabelPoint[right](P){$P  (\alpha : d )$} 
  \tkzDrawPoints[shape=cross](I,J) 
  \tkzLabelPoints(O,I)
  \tkzLabelPoints[left](J) 
\end{tikzpicture}
\end{minipage}%

\subsection{定义命名点：\tkzcname{tkzDefPoint}}

\begin{NewMacroBox}{tkzDefPoint}{\oarg{local options}\parg{$x,y$}\marg{ref} or \parg{$\alpha$:$d$}\marg{ref}}%
\begin{tabular}{lll}%
参数 & 默认值 & 定义  \\
\midrule
\TAline{($x,y$)}{no default}{$x$ 和 $y$ 是两个度量，默认为 cm。}
\TAline{($\alpha$:$d$)}{no default}{$\alpha$ 是以度为单位的角，$d$ 是一个度量}
\TAline{\{ref\}}{no default}{赋值给该点的引用：$A$、$T\_a$、$P1$ 或 $P_1$}
\bottomrule
\end{tabular}

\medskip
\emph{这个宏的必选参数是两个十进制数：第一种情况是两个长度，第二种情况是一个以度为单位的角和一个长度。不要混淆点的引用和名称：引用供计算使用，但名称通常与引用相同。}

\medskip
\begin{tabular}{lll}%
\toprule
选项             & 默认值 & 定义  \\
\midrule
\TOline{label} {no default} {允许你在预定义距离处放置标签}
\TOline{shift} {no default} {将 $(x,y)$ 或 $(\alpha:d)$ 加到所有坐标上}
\end{tabular}
\end{NewMacroBox}

\subsubsection{笛卡尔坐标}
 
\begin{tkzexample}[latex=6cm,small]
  \begin{tikzpicture}
  \tkzInit[xmax=5,ymax=5] % limits the size of the axes
  \tkzDrawX[>=latex]
  \tkzDrawY[>=latex]  
  \tkzDefPoint(0,0){A}
  \tkzDefPoint(4,0){B}
  \tkzDefPoint(0,3){C} 
  \tkzDrawPolygon(A,B,C)
  \tkzDrawPoints(A,B,C)
  \end{tikzpicture}
\end{tkzexample}

\subsubsection{用 \tkzNamePack{xfp} 进行计算}

 \begin{tkzexample}[latex=7cm,small]
\begin{tikzpicture}[scale=1]
  \tkzInit[xmax=4,ymax=4]
  \tkzDrawX\tkzDrawY
  \tkzDefPoint(-1+2,sqrt(4)){O}
  \tkzDefPoint({3*ln(exp(1))},{exp(1)}){A}
  \tkzDefPoint({4*sin(pi/6)},{4*cos(pi/6)}){B}
  \tkzDrawPoints(O,B,A)
\end{tikzpicture}
\end{tkzexample}

\subsubsection{极坐标}

\begin{tkzexample}[latex=7cm,small]
  \begin{tikzpicture}
  \foreach \an [count=\i] in {0,60,...,300}
   { \tkzDefPoint(\an:3){A_\i}}
  \tkzDrawPolygon(A_1,A_...,A_6)
  \tkzDrawPoints(A_1,A_...,A_6)
  \end{tikzpicture}
\end{tkzexample}

\subsubsection{相对点}
首先，我们可以使用 \TIKZ 的 \tkzNameEnv{scope} 环境。
在下面的示例中，我们有一种定义等边三角形的方法。

\begin{tkzexample}[latex=7cm,small]
\begin{tikzpicture}[scale=1]
 \begin{scope}[rotate=30]
  \tkzDefPoint(2,3){A}
  \begin{scope}[shift=(A)]
     \tkzDefPoint(90:5){B}
     \tkzDefPoint(30:5){C}
  \end{scope}
 \end{scope}
 \tkzDrawPolygon(A,B,C)
\tkzLabelPoints[above](B,C)
\tkzLabelPoints[below](A)
\tkzDrawPoints(A,B,C)
\end{tikzpicture}
\end{tkzexample}

\subsection{相对于另一个点的点：\tkzcname{tkzDefShiftPoint}}
\begin{NewMacroBox}{tkzDefShiftPoint}{\oarg{Point}\parg{$x,y$}\marg{ref} or \parg{$\alpha$:$d$}\marg{ref}}%
\begin{tabular}{lll}%
参数 & 默认值 & 定义 \\
\midrule
\TAline{($x,y$)}{no default}{$x$ 和 $y$ 是两个度量，默认为 cm。}
\TAline{($\alpha$:$d$)}{no default}{$\alpha$ 是以度为单位的角，$d$ 是一个度量}
\TAline{\{ref\}}{no default}{赋值给该点的引用：$A$、$T\_a$、$P1$ 或 $P_1$}

\midrule
选项 & 默认值 & 定义 \\

\midrule
\TOline{[pt]} {no default} {\tkzcname{tkzDefShiftPoint}[A](0:4)\{B\}}
\end{tabular}
\end{NewMacroBox}

\subsubsection{等腰三角形}
这个宏将一个点相对于另一个点放置；它等价于一次平移。下面是如何构造一个以 $A$ 为主要顶点、该顶点处角为 $30^{\circ}$ 的等腰三角形。

\begin{tkzexample}[latex=7cm,small]
\begin{tikzpicture}[rotate=-30]
 \tkzDefPoint(2,3){A}
 \tkzDefShiftPoint[A](0:4){B}
 \tkzDefShiftPoint[A](30:4){C}
 \tkzDrawSegments(A,B B,C C,A)
 \tkzMarkSegments[mark=|](A,B A,C)
 \tkzDrawPoints(A,B,C)
 \tkzLabelPoints[right](B,C)
 \tkzLabelPoints[above left](A)
\end{tikzpicture}
\end{tkzexample}

\subsubsection{等边三角形}
让我们看看如何得到一个等边三角形（有更简单的方法，见 \tkzcname{tkzDefTriangle}）。

\begin{tkzexample}[latex=7cm,small]
\begin{tikzpicture}[scale=1]
 \tkzDefPoint(2,3){A}
 \tkzDefShiftPoint[A](30:3){B}
 \tkzDefShiftPoint[A](-30:3){C}
 \tkzDrawPolygon(A,B,C)
 \tkzDrawPoints(A,B,C)
 \tkzLabelPoints[right](B,C)
 \tkzLabelPoints[above left](A)
 \tkzMarkSegments[mark=|](A,B A,C B,C)
\end{tikzpicture}
\end{tkzexample}

\subsubsection{平行四边形}
有更简单的方法，见 \tkzcname{tkzDefParallelogram}。
\begin{tkzexample}[latex=7cm,small]
\begin{tikzpicture}
 \tkzDefPoint(0,0){A}
 \tkzDefPoint(30:3){B}
 \tkzDefShiftPointCoord[B](10:2){C}
 \tkzDefShiftPointCoord[A](10:2){D}
 \tkzDrawPolygon(A,...,D)
 \tkzDrawPoints(A,...,D)
\end{tikzpicture}
\end{tkzexample}

\subsection{多个点的定义：\tkzcname{tkzDefPoints}}

\begin{NewMacroBox}{tkzDefPoints}{\oarg{local options}\marg{$x_1/y_1/r_1,x_2/y_2/r_2$, ...}}%
$x_i$ 和 $y_i$ 是引用为 $r_i$ 的点的坐标。

\begin{tabular}{lll}%
\toprule
参数 & 默认值 & 示例  \\
\midrule
\TAline{$x_i/y_i/r_i$}{}{\tkzcname{tkzDefPoints\{0/0/O,2/2/A\}}}
\end{tabular}

\medskip
\begin{tabular}{lll}%
选项             & 默认值 & 定义   \\
\midrule
\TOline{shift} {no default} {将 $(x,y)$ 或 $(\alpha:d)$ 加到所有坐标上}
\end{tabular}
\end{NewMacroBox}

\subsection{创建三角形}
\begin{tkzexample}[latex=6cm,small]
\begin{tikzpicture}[scale=.75]
 \tkzDefPoints{0/0/A,4/0/B,4/3/C}
 \tkzDrawPolygon(A,B,C)
 \tkzDrawPoints(A,B,C)
\end{tikzpicture}
\end{tkzexample}

\subsection{创建正方形}
请注意这里绘制多边形的语法。
\begin{tkzexample}[latex=6cm,small]
\begin{tikzpicture}[scale=1]
 \tkzDefPoints{0/0/A,2/0/B,2/2/C,0/2/D}
 \tkzDrawPolygon(A,...,D)
 \tkzDrawPoints(A,...,D)
\end{tikzpicture}
\end{tkzexample}

\endinput
